%% Proposed Solution 4.3 - owner: Jivit Rana (2401010202)
%% Scalar baseline, benchmark, cycle-counting rules, prototype status.
%%
%% Phase I has no separate Results section, so the preliminary numbers live
%% here. Everything below is measured, not estimated: reproduce with
%%   make -C sim scalar_top     (scalar cycles)
%%   make -C sim system         (SIMD cycles, checked against the golden model)
%%   python synth/run_synth.py  (area figures)

The value of a SIMD datapath is a claim about throughput, and a claim about
throughput is meaningless without something to measure it against. This section
describes the scalar baseline the vector unit is compared with, the rules under
which both machines are counted, and the figures those rules produce.

\subsubsection{Scalar Baseline}
The baseline (\texttt{scalar\_top.v}) is a single-issue machine that processes
one element at a time through four phases --- load $A[i]$, load $B[i]$, add,
store $C[i]$ --- one cycle per phase. Two design decisions make the comparison
honest rather than flattering:

\begin{itemize}
  \item \textbf{The same ALU.} \texttt{scalar\_top} instantiates the identical
        \texttt{lane\_alu} module used by each of the four vector lanes. The two
        machines therefore differ in \emph{how many} ALUs operate per cycle and
        in nothing else. A separately written scalar adder would have made the
        speed-up a comparison of two different circuits.
  \item \textbf{A 32-bit memory, not 128-bit.} The baseline's data memory is one
        element wide, so a scalar load moves one element per cycle where a
        vector load moves \texttt{LANES} of them. Giving the baseline a wide
        memory would have concealed precisely the effect being measured.
\end{itemize}

Both machines execute on byte-identical data: \texttt{sim/vectors/scalar\_init.hex}
is unpacked from the same golden-model vectors that fill the SIMD machine's
\texttt{dmem}.

\subsubsection{Cycle-Counting Rules}
The rules were fixed \emph{before} the numbers were collected, because rules
chosen after a result is known are rules fitted to that result. They are
recorded in \texttt{docs/cycle\_counting.md} and implemented identically in both
tops.

\begin{table}[htbp]
  \centering
  \caption{Cycle-counting rules, applied identically to both machines}
  \label{tab:cycle_rules}
  \begin{tabular}{@{}lll@{}}
    \toprule
    \textbf{Operation} & \textbf{SIMD} & \textbf{Scalar} \\
    \midrule
    Load        & 1 cycle, 4 elements & 1 cycle, 1 element \\
    ALU         & 1 cycle, 4 lanes    & 1 cycle, 1 element \\
    Store       & 1 cycle, 4 elements & 1 cycle, 1 element \\
    Loop overhead & none              & none \\
    \texttt{HALT}/final cycle & excluded & excluded \\
    \bottomrule
  \end{tabular}
\end{table}

Loop overhead is zero on both sides, and that symmetry is deliberate. The vector
ISA has no branch instruction, so \texttt{bench\_add16} is fully unrolled
straight-line code whose program counter advances alongside each instruction at
no cost. Charging the scalar machine a separate cycle for its index increment
would have billed it for bookkeeping the vector machine does not pay for either,
inflating the speed-up. \texttt{scalar\_top} therefore folds the increment into
its store cycle.

\subsubsection{Measured Results}
Both machines were run under Icarus Verilog~\cite{icarus} on the same
16-element array addition, $C[i] = A[i] + B[i]$.

\begin{table}[htbp]
  \centering
  \caption{Scalar versus SIMD on a 16-element array addition}
  \label{tab:cycle_counts}
  \begin{tabular}{@{}lrrrr@{}}
    \toprule
    \textbf{Machine} & \textbf{Lanes} & \textbf{Cycles} &
    \textbf{Elem/cycle} & \textbf{Speed-up} \\
    \midrule
    \texttt{scalar\_top} & 1 & 64 & 0.25 & 1.00$\times$ \\
    \texttt{simd\_top}   & 4 & \textbf{16} & \textbf{1.00} & \textbf{4.00$\times$} \\
    \bottomrule
  \end{tabular}
\end{table}

The measured speed-up is exactly $4.00\times$, equal to the lane count. Landing
precisely on the theoretical bound invites suspicion, so it is worth stating why
it happens: the benchmark is pure data-parallel work and nothing else. Every
scalar operation has an exact vector counterpart, so the instruction mix is the
same on both sides; loop overhead is zero on both; and the element count is an
exact multiple of \texttt{LANES}, so no lane is ever idle. Change any one of
those and the figure falls --- with 17 elements the vector machine would need a
fifth iteration to process a single element, giving $3.2\times$.

$4.00\times$ should therefore be read as the ceiling being reached under ideal
conditions, not as a general claim that SIMD is four times faster. A workload
with partial vectors, data-dependent control flow, or non-contiguous memory
access would score lower.

\subsubsection{Area Cost}
Throughput is only half of the trade. Synthesising with Yosys~\cite{yosys}
(\texttt{synth\_xilinx}, flattened, no I/O pads) gives the area the speed-up is
bought with.

\begin{table}[htbp]
  \centering
  \caption{Area of one lane against four (Xilinx primitives)}
  \label{tab:area}
  \begin{tabular}{@{}lrrr@{}}
    \toprule
    \textbf{Module} & \textbf{LUTs} & \textbf{CARRY4} & \textbf{Ratio} \\
    \midrule
    \texttt{lane\_alu} & 251 & 19 & 1.00$\times$ \\
    \texttt{vec\_alu} (\texttt{LANES} = 4) & 1012 & 76 &
      \textbf{4.03$\times$ / 4.00$\times$} \\
    \bottomrule
  \end{tabular}
\end{table}

Four lanes cost $4.03\times$ the LUTs and exactly $4\times$ the carry chains of
one. Area scales with the lane count and nothing else appears, which is the
measured evidence for the claim in Section~\ref{sec:datapath} that the
\texttt{generate} loop costs nothing over manual instantiation. Taken together,
a $4.00\times$ throughput gain for a $4.03\times$ area cost is close to a linear
trade on this benchmark.

No timing analysis was performed: this flow targets no device and places and
routes nothing, so no $F_{max}$ is reported. Quoting a frequency the toolchain
never produced would be the least defensible number in this report. Timing
closure on a real device is Phase II work.

\subsubsection{Prototype Status}
The design is complete and verified. All eight testbenches pass with no
failures and no skipped checks; the integrated machine is checked against an
independent Python model of the instruction set
(\texttt{tb\_system}, 18/18 result and cycle-count checks); and a random-program
fuzzer runs generated programs on both \texttt{simd\_top} and that model,
comparing final memory and cycle counts. The fuzzer found a defect the
hand-written testbenches could not: the register file and data memory were being
written while \texttt{rst} was asserted, which stayed hidden because every
hand-written benchmark happens to begin with an idempotent \texttt{VLD}.
